\documentclass[12pt]{article}
\usepackage[a4paper,margin=25mm]{geometry}
\usepackage{ctable,fancyhdr,osf-txfonts}
\parindent0pt
\pagestyle{fancy}
\fancyhead{}\fancyfoot{}
\bibliographystyle{plain}
\begin{document}
This text shows some tests. The heading is used to tell what is being tested, so
use one page per test and use fancyhead to set the heading. 
\clearpage

\fancyhead[C]{notespar sets notes in a block}
\ctable[
  caption=Notes in tabular form,
]{l}{
  \tnote[a]{note.}
  \tnote[b]{other note.}
  \tnote[c]{yet another.}
  \tnote[d]{fourth note; this is a lengthy one, long enough to be wider than the table.}
}{\FL
Some\tmark[a] text\tmark[b] for the first line\NN
and some\tmark[c] for the second line\tmark[d]\LL
}

\ctable[
  caption=Notes in a block,
  notespar
]{l}{
  \tnote[a]{note.}
  \tnote[b]{other note.}
  \tnote[c]{yet another.}
  \tnote[d]{fourth note; this is a lengthy one, long enough to be wider than the table.}
}{\FL
Some\tmark[a] text\tmark[b] for the first line\NN
and some\tmark[c] for the second line\tmark[d]\LL
}

\clearpage
\fancyhead[C]{A figure}
\ctable[figure,botcap,caption=Terminology]{c}{}{\includegraphics{metrics}}

\clearpage

\fancyhead[C]{captionskip; first: without ctable}
\begin{table}[h]\begin{center}
  \caption{x top caption}
  \begin{tabular}{c}\hline a single line tabular\tabularnewline\hline\end{tabular}
  \setcounter{table}{0}
  \caption{x bottom caption}
\end{center}\end{table}

Without ascenders and descenders:
\def\tablename{xxx}
\begin{table}[h]\begin{center}
  \setcounter{table}{0}
  \caption{xxx}
  \begin{tabular}{c}\hline a single line tabular\tabularnewline\hline\end{tabular}
  \setcounter{table}{0}
  \caption{xxx}
\end{center}\end{table}

\begin{table}[h]\begin{center}
  \setcounter{table}{0}
  \caption{xpx}
  \begin{tabular}{c}\hline a single line tabular\tabularnewline\hline\end{tabular}
  \setcounter{table}{0}
  \caption{xTx}
\end{center}\end{table}

\clearpage
\fancyhead[C]{captionskip; now ctable: default 2ex, test 4ex}
\setcounter{table}{0}
\ctable[pos=h!,caption=xxx]{c}{}{\FL a single line ctable\LL}
\setcounter{table}{0}
\ctable[pos=h!,botcap,caption=xxx]{c}{}{\FL a single line ctable\LL}
\setcounter{table}{0}
\ctable[pos=h!,captionskip=4ex,caption=xxx]{c}{}{\FL a single line ctable\LL}
\setcounter{table}{0}
\ctable[pos=h!,botcap,captionskip=4ex,caption=xxx]{c}{}{\FL a single line ctable\LL}

\clearpage
\fancyhead[C]{captionskip; now ctable: default 2ex, test 4ex}
\setcounter{table}{0}
\ctable[pos=h!,caption=xpx]{c}{}{\FL a single line ctable\LL}
\setcounter{table}{0}
\ctable[pos=h!,botcap,caption=xTx]{c}{}{\FL a single line ctable\LL}
\setcounter{table}{0}
\ctable[pos=h!,captionskip=4ex,caption=xpx]{c}{}{\FL a single line ctable\LL}
\setcounter{table}{0}
\ctable[pos=h!,botcap,captionskip=4ex,caption=xTx]{c}{}{\FL a single line ctable\LL}

\clearpage
\fancyhead[C]{CENTERED}

\def\CAP{Zetmeel (enzymatisch), gebruikte enzymen}
\newlength{\capwidth}\settowidth{\capwidth}{Table 3: \CAP}
\ctable[
  pos=h,
  framerule=1pt,
  framebg=1 .9 .9,
  caption={\CAP}
]{llc}{}{\FL
lab& tijd & temperatuur \NN
   & & (C) \ML
02 & 1 & 60 \NN
10 & 4 & 60 \NN
23 & 16 & 60 \NN
26 & 12 & 60 \LL
}

\ctable[
  pos=h,
  framerule=1pt,
  framebg=.9 1 .9,
  mincapwidth=\capwidth,
  caption={\CAP}
]{llc}{}{\FL
lab& tijd & temperatuur \NN
   & & (C) \ML
02 & 1 & 60 \NN
10 & 4 & 60 \NN
23 & 16 & 60 \NN
26 & 12 & 60 \LL
}
\clearpage

\fancyhead[C]{RIGHT ALIGNED}


\ctable[
  pos=h,
  right,
  framerule=1pt,
  framebg=.9 .9 1,
  width=\capwidth,
  caption={\CAP}
]{llX}{}{\FL
lab& tijd & temperatuur \NN
   & & (C) \ML
02 & 1 & 60 \NN
10 & 4 & 60 \NN
23 & 16 & 60 \NN
26 & 12 & 60 \LL
}

\ctable[
  pos=h,
  right,
  framerule=1pt,
  framebg=1 .9 .9,
  caption={\CAP}
]{llc}{}{\FL
lab& tijd & temperatuur \NN
   & & (C) \ML
02 & 1 & 60 \NN
10 & 4 & 60 \NN
23 & 16 & 60 \NN
26 & 12 & 60 \LL
}

\ctable[
  pos=h,
  right,
  framerule=1pt,
  framebg=.9 1 .9,
  mincapwidth=\capwidth,
  caption={\CAP}
]{llc}{}{\FL
lab& tijd & temperatuur \NN
   & & (C) \ML
02 & 1 & 60 \NN
10 & 4 & 60 \NN
23 & 16 & 60 \NN
26 & 12 & 60 \LL
}

\clearpage
\fancyhead[C]{LEFT ALIGNED}


\ctable[
  pos=h,
  left,
  framerule=1pt,
  framebg=.9 .9 1,
  width=\capwidth,
  caption={\CAP}
]{llX}{}{\FL
lab& tijd & temperatuur \NN
   & & (C) \ML
02 & 1 & 60 \NN
10 & 4 & 60 \NN
23 & 16 & 60 \NN
26 & 12 & 60 \LL
}

\ctable[
  pos=h,
  left,
  framerule=1pt,
  framebg=1 .9 .9,
  caption={\CAP}
]{llc}{}{\FL
lab& tijd & temperatuur \NN
   & & (C) \ML
02 & 1 & 60 \NN
10 & 4 & 60 \NN
23 & 16 & 60 \NN
26 & 12 & 60 \LL
}

\ctable[
  pos=h,
  left,
  framerule=1pt,
  framebg=.9 1 .9,
  mincapwidth=\capwidth,
  caption={\CAP}
]{llc}{}{\FL
lab& tijd & temperatuur \NN
   & & (C) \ML
02 & 1 & 60 \NN
10 & 4 & 60 \NN
23 & 16 & 60 \NN
26 & 12 & 60 \LL
}

\clearpage

\fancyhead[C]{BIG CAPTION AND FOOTNOTE, SMALL TABLE, MINCAPWIDTH=111mm}


 
\def\CAP{%
We thrive in information||thick worlds because of our
marvelous and everyday capacity to select, edit,
single out, structure, highlight, group, pair, merge,
harmonize, synthesize, focus, organize, condense,
reduce, boil down, choose, categorize, catalog, classify,
list, abstract, scan, look into, idealize, isolate,
discriminate, distinguish, screen, pigeonhole, pick over,
sort, integrate, blend, inspect, filter, lump, skip,
smooth, chunk, average, approximate, cluster, aggregate,
outline, summarize, itemize, review, dip into,
flip through, browse, glance into, leaf through, skim,
refine, enumerate, glean, synopsize, winnow the wheat
from the chaff and separate the sheep from the goats.}
{\LARGE
\ctable[
  caption={\CAP},
  mincapwidth=111mm
]{l} {\tnote{\input{tufte}}}
{\FL
Anzahl\NN
\%\LL
}
}
\end{document}
$Id$
